% Generated draft: requires mathematical and source review.
\documentclass[11pt]{article}
\usepackage[margin=25mm]{geometry}
\usepackage{fontspec}
\setmainfont{Latin Modern Roman}
\usepackage{amsmath,amssymb,mathtools,unicode-math}
\setmathfont{Latin Modern Math}
\usepackage{enumitem,needspace,xurl,hyperref}
\setlength{\parindent}{0pt}
\setlength{\parskip}{5pt}
\setlength{\emergencystretch}{3em}
\newcommand{\N}{\mathbb{N}}
\newcommand{\Z}{\mathbb{Z}}
\newcommand{\Q}{\mathbb{Q}}
\newcommand{\R}{\mathbb{R}}
\newcommand{\C}{\mathbb{C}}
\newcommand{\F}{\mathbb{F}}
\newcommand{\A}{\mathbb{A}}
\newcommand{\PP}{\mathbb{P}}
\newcommand{\E}{\mathbb{E}}
\newcommand{\eps}{\varepsilon}
\newcommand{\dd}{\,\mathrm{d}}
\begin{document}
\section*{UnsolvedMath ID: 30006408}
\textbf{Source problem code:} catalogue entry
\subsection{Definitions and mathematical statement}
For a specified Riemannian manifold \((M,g)\) and Euclidean target \(\R^N\), an isometric immersion is a \(C^1\) immersion \(u:M\to\R^N\) satisfying
\[
\langle \dd u_x(v),\dd u_x(w)\rangle_{\R^N}=g_x(v,w)
\qquad(x\in M,\ v,w\in T_xM).
\]
For \(0<\alpha<1\), the notation \(C^{1,\alpha}\) means that the first derivatives are locally H\"older continuous with exponent \(\alpha\) in coordinate charts. This is the usual Euclidean-target interpretation of the terminology; the supplied passage does not specify the domain, target, or global regularity requirements.

The source asks whether there exists a critical H\"older exponent \(\alpha_0\) separating rigidity from flexibility. Its proposed dichotomy can be written schematically as
\[
\exists\alpha_0\quad\forall\alpha\in(0,1):\qquad
\begin{cases}
\alpha>\alpha_0\ \Longrightarrow\ \text{rigidity for the intended }C^{1,\alpha}\text{ isometric immersions},\\
\alpha<\alpha_0\ \Longrightarrow\ \text{a }C^{1,\alpha}\text{ extension of the Nash--Kuiper theorem}.
\end{cases}
\]
The source explicitly qualifies the first conclusion as ``some form of'' rigidity. Neither that property nor the precise flexibility theorem is defined in the supplied record. In the usual approximation formulation, Nash--Kuiper flexibility concerns uniform approximation of short maps by isometric immersions; a map \(v:M\to\R^N\) is strictly short when
\[
\langle \dd v_x(\xi),\dd v_x(\xi)\rangle_{\R^N}<g_x(\xi,\xi)
\qquad(x\in M,\ 0\ne\xi\in T_xM).
\]
This explanation does not specify which initial maps, approximation topology, boundary conditions, or dimension restrictions the source intends.

Consequently, the displayed dichotomy is an unresolved formulation of a question, not a fully specified universally quantified conjecture. The admissible geometric class, metric regularity, dimensions and codimension, curvature or other geometric hypotheses, meaning of rigidity, and precise flexibility conclusion must be supplied before such a conjecture can be stated. The record does not determine whether the threshold is universal or depends on the geometric class, nor does it specify its admissible range or any assertion at \(\alpha=\alpha_0\). Its motivating contrast between flexibility in \(C^1\) and rigidity in \(C^2\) must likewise be understood subject to unspecified geometric hypotheses, rather than as an assertion about all isometric immersions.
\subsection{Short English statement}
The question asks whether a critical H\"older exponent separates rigidity of isometric immersions above the threshold from Nash--Kuiper flexibility below it. The supplied passage leaves the geometric setting and the meaning of rigidity unspecified, so it does not determine a single precise conjecture. The supplied status is open; this draft does not independently verify that status.
\subsection{Sources}
\begin{itemize}
\item Supplied local UnsolvedMath registry record, permanent ID \(30006408\), title: The H\"older Threshold for Isometric-Immersion Rigidity. Exact source problem code: \texttt{catalogue entry}. Collection: Oberwolfach Reports Open Problems. The record gives proposed year 2025 and names no individual proposer.
\item Supplied report citation: \emph{Partial Differential Equations (2026)}, Oberwolfach Reports, DOI \texttt{10.4171/owr/2025/33}, pp.\ 1753--1822. \url{https://doi.org/10.4171/owr/2025/33}. The record attributes the question to workshop-report discussion and describes its original passage as retained verbatim. No exact page for the passage is supplied. The citation's 2026 designation and the 2025 DOI/proposed-year information are retained without reconciliation.
\item Supplied literature reference: \emph{Flexibility of Codimension One C one theta Isometric Immersions}. \url{https://arxiv.org/abs/2603.08382}.
\item Supplied literature reference: \emph{A Nash-Kuiper theorem for isometric immersions beyond Borisov's exponent}. \url{https://arxiv.org/abs/2503.13867}.
\item Literature assessment dated exactly \texttt{2026-08-22}: the supplied curation labels the problem open, states that the optimal \(C^{1,\alpha}\) flexibility--rigidity threshold is unknown, and reports improvements in flexible exponents. This is supplied dated curation derived from the OWR record, not an independently verified literature assessment. The record itself requests independent verification of time-sensitive status and cited claims.
\item The supplied statement status is \texttt{corrected\_verified}; its verification note says that the official report was checked and the question isolated with repaired notation. These are supplied provenance claims, not verification performed for this draft.
\end{itemize}
\end{document}
